\begin{examplebox}
	\textbf{\textcolor{CExample}{例1（基础）}}

	\medskip
	\textbf{\textcolor{CExample}{题目：}}
	在三棱锥$P-ABC$中，已知$PA\perp AB$，$PA\perp AC$，$AB\cap AC=A$，$AB\subset$平面$ABC$，$AC\subset$平面$ABC$。求证：$PA\perp$平面$ABC$。

	\medskip
	\textbf{\textcolor{CExample}{解题步骤：}}
	\begin{description}[style=nextline, leftmargin=2.8em, labelsep=0.5em, itemsep=0.45em]
		\item[\textbf{第一步（找相交直线）：}]
			在平面$ABC$内找到两条相交直线$AB$和$AC$，且$AB\cap AC=A$。
			\par\textbf{理由：} 线面垂直判定定理要求平面内有两条相交直线与已知直线垂直。
		\item[\textbf{第二步（验证垂直条件）：}]
			由已知条件直接得到$PA\perp AB$，$PA\perp AC$。
			\par\textbf{理由：} 题目已明确给出垂直关系。
		\item[\textbf{第三步（使用定理）：}]
			因为$PA$垂直于平面$ABC$内的两条相交直线$AB$和$AC$，根据线面垂直判定定理，$PA\perp$平面$ABC$。
			\par\textbf{理由：} 判定定理：“若一条直线同时垂直于平面内的两条相交直线，则该直线垂直于平面”。
	\end{description}

	\medskip
	\textbf{\textcolor{CExample}{关键结论：}}
	$PA\perp$平面$ABC$。

	\medskip
	\textbf{\textcolor{CExample}{例2（稍变形）}}

	\medskip
	\textbf{\textcolor{CExample}{题目：}}
	在四棱锥$P-ABCD$中，底面$ABCD$是矩形，$PA\perp AB$，$PA\perp AD$。求证：$PA\perp$平面$ABCD$。

	\medskip
	\textbf{\textcolor{CExample}{解题步骤：}}
	\begin{description}[style=nextline, leftmargin=2.8em, labelsep=0.5em, itemsep=0.45em]
		\item[\textbf{第一步（找相交直线）：}]
			在底面$ABCD$内，$AB$和$AD$相交于点$A$，且$AB\subset$平面$ABCD$，$AD\subset$平面$ABCD$。
			\par\textbf{理由：} 矩形相邻两边相交于顶点，满足“相交直线”条件。
		\item[\textbf{第二步（确认垂直）：}]
			已知$PA\perp AB$，$PA\perp AD$。
			\par\textbf{理由：} 题目条件直接给出。
		\item[\textbf{第三步（应用定理）：}]
			由于$PA$垂直于平面$ABCD$内的两条相交直线$AB$和$AD$，由判定定理得$PA\perp$平面$ABCD$。
			\par\textbf{理由：} 线面垂直判定定理直接应用。
	\end{description}

	\medskip
	\textbf{\textcolor{CExample}{关键结论：}}
	$PA\perp$平面$ABCD$。

	\medskip
	\textbf{\textcolor{CExample}{例3（综合应用）}}

	\medskip
	\textbf{\textcolor{CExample}{题目：}}
	在三棱锥$P-ABC$中，$\angle ABC=90^\circ$，$PA\perp$平面$ABC$。求证：$BC\perp$平面$PAB$。

	\medskip
	\textbf{\textcolor{CExample}{解题步骤：}}
	\begin{description}[style=nextline, leftmargin=2.8em, labelsep=0.5em, itemsep=0.45em]
		\item[\textbf{第一步（由性质得线线垂直）：}]
			因为$PA\perp$平面$ABC$，所以$PA\perp BC$。
			\par\textbf{理由：} 若一条直线垂直于一个平面，则它垂直于平面内的任意一条直线。
		\item[\textbf{第二步（利用已知垂直）：}]
			又$\angle ABC=90^\circ$，即$AB\perp BC$。
			\par\textbf{理由：} 直角三角形的直角边互相垂直。
		\item[\textbf{第三步（判定定理）：}]
			$BC$垂直于平面$PAB$内的两条相交直线$PA$和$AB$（$PA\cap AB=A$），根据判定定理，$BC\perp$平面$PAB$。
			\par\textbf{理由：} 线面垂直判定定理。
	\end{description}

	\medskip
	\textbf{\textcolor{CExample}{关键结论：}}
	$BC\perp$平面$PAB$。
\end{examplebox}
